% ============================================================
% DRAFT — Synthesis paragraphs for cap3_SLR.tex
% Generated 2026-04-19 from extraction_template + top30 ranking
% Insert into \subsection{RQ2.2} and \subsection{RQ3.1/RQ3.2}
% ============================================================

% -----------------------------------------------------------
% ANCHOR PARAGRAPH — seminal arc (insert at start of RQ2.1 or
% as opening of Results/Synthesis section)
% -----------------------------------------------------------

The study of software processes through computational methods dates to the
mid-1990s. Cook and Wolf~\cite{cook1995,cook1998} were the first to
demonstrate that process models could be \emph{discovered automatically} from
event traces---using finite-state automata, Markov models, and neural
networks---establishing the theoretical vocabulary that later became process
mining. Rubin et al.~\cite{rubin2007} generalised this foundation into the
first explicitly named \emph{process mining framework for software processes},
applying the ProM toolkit to software configuration management logs and
showing that alpha-algorithm-based discovery, heuristic mining, and social
network analysis could all be adapted to SDLC event logs.
These three works collectively define the starting point of the literature
and are treated here as the canonical origin of the field.


% -----------------------------------------------------------
% RQ2.2 — Stochastic techniques (IC2+IC3 cluster, 14 papers)
% INSERT into \subsubsection{RQ2.2: Predictive and Stochastic Techniques}
% replaces / expands the current paragraph that starts
% "The PDF analysis sharpens the stochastic-method picture."
% -----------------------------------------------------------

The stochastic-technique picture that emerges from the top-ranked literature
divides naturally into three sub-families. The largest is
\textbf{Markov-chain-based prediction}, which spans from defect-removal
modelling~\cite{tian2005,jeon2015} through agile risk
assessment~\cite{lunesu2021} and CI/CD reliability
analysis~\cite{bhadra2022,bhadra2023} to reinforcement-learning formulations
of pull-request outcome prediction~\cite{joshi2024}. Common to all of these
is the use of discrete-state transition matrices estimated from repository or
defect-tracking data; the prediction target varies from remaining defect count
and release date to the probability of PR acceptance within a sprint
boundary. The second sub-family uses \textbf{stochastic Petri
nets}~\cite{benmesmia2021,bhadra2022,bhadra2023}: these models capture
parallelism and resource contention in CI/CD and DevOps workflows more
faithfully than plain Markov chains, but at the cost of additional
specification effort and sensitivity to parametric assumptions. The third
sub-family employs \textbf{Monte Carlo simulation} directly on project
data~\cite{magennis2015,lunesu2021,li2020}, typically feeding historical cycle
times into throughput forecasts. Ghosh~\cite{ghosh2009} and Santos et
al.~\cite{santos2018} extend this line to globally distributed projects and
Waterfall-style estimation, respectively, while Li et al.~\cite{li2020}
demonstrate the construction of a \emph{hybrid} simulation combining
System Dynamics and Discrete Event Simulation for industrial project
forecasting---the closest prior work to PATHCAST's Stages 3--4 in terms of
modelling ambition.

A consistent limitation across all three sub-families is their
\emph{disconnection from process mining}: the stochastic models are
parametrised from aggregate statistics (average cycle times, historical
defect counts) rather than from mined transition structures. The consequence
is that model states correspond to coarse lifecycle phases rather than to
discovered process states with empirical fitness guarantees. This disconnect
is the operational definition of Finding~F3 and justifies the absorbing-chain
construction in Stage~2 of PATHCAST, where state-transition probabilities are
derived directly from the mined process model rather than estimated
independently.


% -----------------------------------------------------------
% RQ2.2 continued — PM+stochastic integration (IC1+IC2 cluster, 4+1 papers)
% -----------------------------------------------------------

A smaller but highly relevant cluster combines process mining and stochastic
analysis within the same study. Incerto et al.~\cite{incerto2025} propose
stochastic conformance checking based on variable-length Markov chains
(VLMCs), outperforming 18 state-of-the-art techniques on 10 of 11 benchmark
datasets---arguably the most technically advanced integration of PM and
Markov modelling in the current literature, though it targets conformance
assessment rather than operational forecasting. L\'opez-Pintado and
Dumas~\cite{lopezpintado2023} bridge PM and simulation by discovering
probabilistic resource-availability calendars from event logs, showing
that simulation models with stochastic calendars replicate temporal
distributions of activity instances more faithfully than crisp-calendar
alternatives. Jalote and Kamma~\cite{jalote2021} apply Markov chains to
model task processes extracted from IDE interaction logs, providing empirical
evidence that process similarity within productivity groups is quantifiable
via chain distance metrics.

The single paper that approaches L3 integration is PRIMAD~\cite{guinea2025},
which proposes a Digital Twin for the SDLC by composing pm4py-based process
mining, Akka actor simulation, and LightGBM prediction. PRIMAD is the closest
prior work to PATHCAST in architectural scope; however, it remains at the
component-validation stage (each layer is evaluated independently) and does
not establish formal inter-stage data contracts or a Monte Carlo forecasting
layer. This positions PRIMAD at L2 in the SPMF taxonomy and leaves the L3
slot---formally specified, end-to-end, multi-layer pipeline with empirical
validation---unoccupied in the literature.


% -----------------------------------------------------------
% RQ2.2 continued — PM+forecasting (IC1+IC3 cluster, 5 papers)
% -----------------------------------------------------------

Studies that combine process mining with forecasting objectives (IC1 $\cap$
IC3) tend to approach prediction through \emph{process enhancement} rather
than through explicit stochastic modelling. Gupta~\cite{gupta2017} applies
inductive miner discovery with conformance checking on multi-repository data
from VCS and issue trackers, then feeds process metrics into a predictive
analytics layer for software maintenance---one of the few studies to chain
discovery and prediction in the SE domain, though without a probabilistic
forecasting component. Pourbafrani et al.~\cite{pourbafrani2021} present
SIMPT, a tool that extracts process trees with time parameters from event
logs and enables interactive what-if simulation; its successor
work~\cite{buliga2025} further integrates generative AI with business process
simulation to produce conformance-preserving what-if scenarios. Caldeira et
al.~\cite{caldeira2022} mine IDE interaction logs and show that process
complexity metrics derived from ProM correlate moderately ($r \approx 0.43$)
with software cyclomatic complexity, demonstrating the feasibility of
process-derived features for product-quality prediction.

Taken together, this cluster confirms that PM and prediction are being
connected in the SE literature, but always at the level of sequential
application (PM output feeds a downstream model) rather than through a
unified pipeline with formal contracts. The absence of an absorbing-state
formulation means that none of these studies can express ``probability of
completion within $k$ cycles'' as a first-class model output.


% -----------------------------------------------------------
% RQ1.2 — Event log sources (IC2+IC4 cluster, 3 papers)
% INSERT into \subsubsection{RQ1.2: Event Log Sources and Construction}
% -----------------------------------------------------------

Three high-ranked studies use GitHub repository data to construct Markov
models directly from commit and collaboration traces, without an explicit
process mining step. Jo et al.~\cite{jo2023,jo2024} model GitHub activity
sequences as discrete-time Markov chains and apply probabilistic Computation
Tree Logic for model checking, producing insights about long-run activity
probabilities and project-level transition dynamics across five repositories.
Ortu et al.~\cite{ortu2023} analyse fault-insertion and fault-fixing patterns
in Apache Software Foundation projects through Markov chains and social
network analysis, finding that fault-inserting developers are less likely to
also fix their faults as repository size grows. These studies demonstrate that
SDLC event logs extracted from public repositories are rich enough to sustain
state-transition modelling---a prerequisite for Stage~1 of PATHCAST---while
also revealing that the construction step (deciding which GitHub events
constitute process states) is itself a methodological contribution that the
literature has not yet standardised.


% -----------------------------------------------------------
% RQ3.1 — Integration level table note (for \subsection{RQ3.1})
% Replaces / extends current paragraph on integration levels
% -----------------------------------------------------------

Mapping the top-ranked papers onto the L0--L3 scale confirms and sharpens
the picture drawn from the broader 73-paper PDF subset.
The IC2+IC3 cluster (14 papers) concentrates almost exclusively at L1: each
study applies one stochastic method to software data but does not compose it
with a process mining layer. The IC1+IC2 cluster (4 papers) contains the
two highest-integration examples in the full set---Incerto et
al.~\cite{incerto2025} and L\'opez-Pintado and Dumas~\cite{lopezpintado2023}---both
reaching L2 but stopping short of a complete SDLC forecasting pipeline.
The IC1+IC3 cluster illustrates the complementary limitation: process mining
is used to describe and assess software processes, but the stochastic layer
required to convert discovered models into distributional forecasts is absent.
PRIMAD~\cite{guinea2025} is the single study that assembles elements from
all three ICs, yet its evaluation is component-wise rather than end-to-end.
The practical consequence is unambiguous: the space of L3 systems---unified
PM $\rightarrow$ Markov $\rightarrow$ MC $\rightarrow$ ML architectures
validated on SDLC data---is empty in the current literature, and PATHCAST
is proposed as its first inhabitant.


% -----------------------------------------------------------
% RQ3.2 — Gap table (structured summary, for gap discussion)
% Can be turned into a \begin{table}...\end{table} in LaTeX
% -----------------------------------------------------------

% Gap | Evidence from top-30 | PATHCAST response
% ----+----------------------+------------------
% G1  | 14 IC2+IC3 papers use stochastic models parametrised from
%     | aggregate statistics, not from mined transitions. No paper
%     | derives Markov matrices from a discovered process model.
%     | → PATHCAST Stage 2: absorbing chain from mined net.
%
% G2  | IC1+IC2 papers (Incerto, López-Pintado) achieve stochastic
%     | PM integration but target conformance/simulation, not
%     | multi-step distributional forecasting.
%     | → PATHCAST Stage 3: Monte Carlo over absorbing chain.
%
% G3  | IC1+IC3 papers (Gupta, Caldeira, Pourbafrani) chain PM
%     | and prediction sequentially but without a stochastic layer.
%     | → PATHCAST Stage 4: ML residual correction on MC output.
%
% G4  | No paper uses process-derived features (conformance score,
%     | rework count, Markov expected remaining time, path entropy)
%     | as ML inputs. Caldeira (2022) is closest (process complexity
%     | metrics), but stops at correlation analysis.
%     | → PATHCAST Feature Engineering component.
%
% G5  | Event-log construction from heterogeneous SDLC sources is
%     | treated ad hoc in all reviewed papers; no reusable framework.
%     | → PATHCAST Stage 1: explicit transformation pipeline.
